\subsubsection{算法对比实验设计}

为支撑报告中的 P1 横向比较，本文新增脚本 \verb|run_p1_experiments.py|，位于 \verb|src/p1_candidate_paths/|。该脚本只读取已有的 P1 处理结果，不覆盖任何原始结果；所有新增实验产物写入 \verb|outputs/p1_experiments/|。输出包括算法总体指标表、代表性 OD 对路径对比图和对应 CSV。

\begin{table}[H]
\centering
\small
\begin{tabular}{p{0.18\linewidth}rrrrr}
\toprule
算法 & 路径数 & 平均距离 & 平均转角数 & 平均动态代价 & 平均总成本 \\
\midrule
Grid A* & 272 & 6.463 & 9.206 & 0.000 & 6.463 \\
VG & 272 & 6.054 & 1.846 & 0.000 & 6.054 \\
AVG & 272 & 6.071 & 1.765 & 0.000 & 6.695 \\
DAVG & 272 & 6.071 & 1.765 & 0.000 & 6.695 \\
\bottomrule
\end{tabular}
\caption{P1 四类单算法总体指标对比}
\label{tab:p1-algorithm-summary}
\end{table}

表 \ref{tab:p1-algorithm-summary} 体现了不同算法的主要差异。Grid A* 的平均距离为 6.463，高于 VG 的 6.054；平均转角数为 9.206，也显著高于 VG 的 1.846。这说明栅格离散会引入阶梯状路径和更多方向变化。VG 在连续可视图上寻找几何短路，所以距离和折点都更少。AVG 的平均距离略高于 VG，但平均转角数进一步降低到 1.765，说明转角惩罚会让搜索在少量距离损失下获得更平滑的路径。DAVG 在当前静态算例中与 AVG 一致，因为动态风险代价没有被触发。

\begin{figure}[!htbp]
\centering
\includegraphics[width=0.88\linewidth]{p1_algorithm_metric_summary.png}
\caption{P1 四类算法在距离、转角、动态代价和总成本上的总体对比}
\label{fig:p1-algorithm-metrics}
\end{figure}

图 \ref{fig:p1-algorithm-metrics} 从可视化角度进一步说明：如果只看距离，VG 是当前算例中最短的单路径源；如果考虑转向平滑性，AVG/DAVG 比 VG 稍微牺牲距离，但用转角惩罚表达执行偏好；如果考虑动态风险，当前样本中四类路径动态代价均为 0，说明风险区域并未影响这些静态 OD 最优路径。这为比较不同路径源的调度影响提供了基础。

\subsubsection{路径生成结果概述}

P1 的实验重点是验证二维仓库场景中的候选路径生成是否稳定、完整且可复现。本文分别使用 A*、VG、AVG 和 DAVG 生成关键节点之间的基础路径源，并输出路径成本表与轨迹采样表。不同算法在路径长度、转弯数量和风险规避倾向上存在差异，因此能够用于比较不同路径生成机制的几何特征和代价结构。

需要说明的是，mixed 多候选路径库的融合、路径源对调度结果的影响，以及不同路径源下 $C_{\max}$、$F_2$、$F_3$ 的比较，均属于任务级调度实验，本文在 P2 实验章节中统一展开。P1 在此处只说明基础路径源的生成质量、算法差异和输出文件完整性。

